\documentclass[11pt]{article}
\usepackage[margin=1in]{geometry}
\usepackage{amsmath,amssymb,amsthm}
\usepackage{booktabs}
\usepackage{hyperref}
\usepackage[numbers,sort&compress]{natbib}
\hypersetup{hidelinks}
\emergencystretch=3em

\newtheorem{remark}{Remark}

\newcommand{\Hyp}{\Omega}
\newcommand{\Feas}{\mathcal N}
\newcommand{\Must}{\mathsf{Must}}
\newcommand{\May}{\mathsf{May}}
\newcommand{\Cannot}{\mathsf{Cannot}}
\newcommand{\Uncertain}{\mathsf{U}}

\title{Shrinking the Feasible Family:\\
Answer-Preserving Reduction Techniques for Set-Theoretic Estimation\\
\large A methods note, grounded in the Lean~4 mechanization of project
\texttt{thejeb}}
\author{Project \texttt{thejeb}}
\date{July 28, 2026}

\begin{document}
\maketitle

\begin{abstract}
In set-theoretic estimation (STE) a problem is a family of property sets
$S_i$ over a solution space, and its answer object is the feasibility set
$\Feas=\bigcap_i S_i$ --- the family of hypotheses consistent with every
piece of information \citep{combettes1993}.  Generating that family is
step one of an STE attack.  This note addresses the step that follows:
once the (possibly enormous) feasible family exists, how does one
shrink, restrict, collapse, certify, and structurally compress it while
provably preserving the STE answers --- feasibility, and the must / may
/ cannot queries quantified over feasible hypotheses?  We synthesize the
project's machine-checked toolkit into a reusable technique: exact
property-set application (the STE meet), contravariant restriction
across corpus states, garbage collection of the variable universe, the
must-quotient as the only generally safe deterministic collapse, minimal
supports as shrunken certificates, gluing of local families with an
exact cross-constraint obstruction, and --- past the proven $2^N$
extensional-state wall, where generic reduction provably stops ---
width-bounded structural compression by variable elimination, junction
trees, and factorization, down to Carlson's substitution reduction of
cipher classes \citep{carlson2012}.  Every guarantee is cited to a named
Lean~4 theorem in the project source; qualifications (consistency
premises, the compatibility notion for gluing, the exact scope of the
lower bounds) are stated where the mechanization states them.
\end{abstract}

\section{Setup: the feasible family and the queries to preserve}
\label{sec:setup}

An STE problem over a solution space $\Hyp$ is an indexed family of
property sets $S:I\to\mathcal P(\Hyp)$, one per piece of information;
the feasibility set is
\[
\Feas \;=\; \bigcap_{i\in I} S_i,
\]
and a set-theoretic estimate is any point of $\Feas$
\citep[\S II-C]{combettes1993}.  In the mechanization this is
\path{feasibilitySet} with the membership law
\path{mem_feasibilitySet}: a hypothesis is feasible iff it lies in every
property set.  Enforcing only a subfamily $J\subseteq I$ yields the
partial feasibility set \path{partialFeasibilitySet}, and information
monotonicity --- more information never enlarges the feasible family ---
is the antitonicity theorem \path{partialFeasibilitySet_antitone}.  If
the full intersection is empty, at least one piece of information is
unfair for every candidate truth
(\path{exists_unfair_of_feasibilitySet_eq_empty}): inconsistency is not
a failure mode but a detection principle.

The feasible family is rarely wanted for its own sake.  What an STE
consumer asks of it are \emph{queries}:

\begin{itemize}
\item \textbf{Feasibility}: is $\Feas\neq\emptyset$ (the formulation is
  consistent)?
\item \textbf{Must}: does a relation hold in \emph{every} feasible
  hypothesis?
\item \textbf{May}: does it hold in \emph{some} feasible hypothesis?
\item \textbf{Cannot}: does it fail in every feasible hypothesis?
\end{itemize}

The project's working instantiation is the dynamic-frame model
\path{STE.DynamicFrame.Model}: documents carry claims, constraints carry
provenance supports, a corpus $D$ activates exactly the constraints
whose support lies in $D$, and the feasible family
\path{Model.feasible} is the partial feasibility set of the active
property sets.  Each feasible hypothesis carries an equivalence relation
on claims, and the queries become \path{MustSame}, \path{MaySame},
\path{CannotSame}, and the explicit two-sided marker \path{Uncertain}
(both merge and split remain feasible).  Everything below is stated for
this model but the technique is generic: any STE problem whose
hypotheses support a relational query has the same shape.

The goal of this note is \emph{answer-preserving reduction}: replace the
feasible family, or its representation, by something smaller, while
provably returning the same answers to the queries above.  Reductions
come in three kinds --- reductions that shrink the family itself
(Sections~\ref{sec:meet}--\ref{sec:restriction}), reductions that
compress a query-equivalent summary of it
(Sections~\ref{sec:quotient}--\ref{sec:gluing}), and reductions that
compress its representation
(Sections~\ref{sec:wall}--\ref{sec:structure}).

\section{Base reduction: the exact meet}
\label{sec:meet}

The primitive shrinking step is property-set application: intersect the
current family $X$ with the property set of one constraint $k$,
\[
X \mapsto X \cap S_k .
\]
This is \path{applyConstraint} in \path{Ste.DynamicFrame.Solver}, and
the mechanization proves that this single operation already has every
algebraic property a practitioner needs to schedule it freely:

\begin{itemize}
\item \textbf{Monotone shrinking.}  Application only removes hypotheses
  (\path{applyConstraint_subset}); at the corpus level, activating more
  documents only shrinks the feasible family
  (\path{feasible_antitone}).
\item \textbf{Idempotence.}  Applying the same property set twice equals
  applying it once (\path{applyConstraint_idempotent}): re-running an
  audited constraint is free.
\item \textbf{Order-independence.}  Applications commute
  (\path{applyConstraint_comm}); a finite run depends only on the
  \emph{set} of constraints applied, not their order or multiplicity
  (\path{solve_congr}, \path{solve_perm}).
\item \textbf{Exactness.}  A finite run is extensionally the partial
  feasibility set of the applied constraints
  (\path{solve_eq_partialFeasibilitySet}), and if the schedule
  enumerates exactly the active constraints of corpus $D$, the solver
  output \emph{is} the model's feasible family
  (\path{solve_eq_feasible}).  There is no approximation layer between
  ``run the solver'' and ``the STE answer''.
\item \textbf{Delta decomposition.}  Each application splits the state
  exactly into survivors and casualties: \path{apply_union_reduction}
  proves $X = (X\cap S_k)\;\cup\;(X\setminus(X\cap S_k))$, and
  \path{mem_solverReduction} identifies the discarded set as precisely
  the current candidates violating the new constraint.  A constraint
  causes no reduction iff the current family already refines it
  (\path{solverReduction_eq_empty_iff}) --- the operational test for a
  redundant constraint.
\end{itemize}

\emph{When to reach for it.}  Always.  This is the sound kernel of every
other technique: any pipeline stage that cannot be expressed as (or
proven equal to) a meet with an audited property set is where soundness
must be argued separately.  The delta decomposition additionally gives
exact incremental accounting --- what each constraint killed --- for
free.

\section{Restriction across instances: garbage-collecting the universe}
\label{sec:restriction}

An STE attack rarely runs on one fixed instance.  Corpora grow and
shrink; the practitioner wants to reuse a computed feasible family
rather than regenerate it.  The mechanization gives this a precise
shape: for corpora $D\subseteq E$, the map that views an $E$-feasible
hypothesis as a $D$-feasible one (\path{restrictFeasible}) satisfies the
identity and composition laws \path{restrictFeasible_id} and
\path{restrictFeasible_comp} --- so the assignment
$D\mapsto\mathrm{feasible}(D)$ is a contravariant, presheaf-shaped
construction on the poset of corpora, with the same laws holding for the
packaged world views (\path{restrictWorld_id},
\path{restrictWorld_comp}).  The image of restriction is characterized
exactly: a small-corpus hypothesis is in the image iff the same ambient
hypothesis is feasible on the larger corpus
(\path{mem_range_restrictFeasible_iff}).

Two consequences matter for shrinking:

\textbf{Incremental deltas instead of regeneration.}  Growing the corpus
decomposes exactly: a hypothesis is feasible on $E\supseteq D$ iff it is
feasible on $D$ and satisfies each \emph{newly activated} constraint
(\path{mem_feasible_extension}, with the new constraints isolated as
\path{newConstraints}).  If an insertion activates nothing new, the
family is unchanged (\path{feasible_eq_of_no_new_constraints}).  So the
already-shrunk family is the correct starting point for the meet of
Section~\ref{sec:meet}, never a lossy cache.

\textbf{Partial restriction of the variable universe.}  The claim
universe itself can be garbage-collected.  Restricting active claims
under corpus deletion is deliberately \emph{partial}: a claim restricts
iff its source document remains active, and when the restriction exists
it is unique (\path{claimRestriction_exists_iff},
\path{claimRestriction_unique} --- restriction fibers have size $0$ or
$1$).  Dropping the claims of deleted documents loses no explanatory
power over the surviving universe, because local coreference is
preserved definitionally under world restriction and claim inclusion
(\path{localSame_restrict}).

\begin{remark}[Retraction asymmetry]
Restriction runs against a one-way door.  Adding documents induces a
covariant map on the safe quotients (\path{canonicalMap}, built from
\path{includeActiveClaim} and the monotonicity \path{mustSame_mono});
that map may identify old classes and is not generally injective.  In
the deletion direction there is in general \emph{no} quotient map at
all: a pair forced together on $E$ can become uncertain on
$D\subset E$, and mapping an $E$-class backward would require splitting
it, which is not a function on classes.  Deletion therefore cannot be
implemented as the inverse of a merge history; it must be implemented
semantically, by re-restricting the feasible family --- which the laws
above make exact.
\end{remark}

\emph{When to reach for it.}  Whenever instances are related by
inclusion: streaming corpora, leave-one-out audits, provenance
ablations.  Restrict, do not recompute; garbage-collect claims of
retired documents; and never try to undo a merge by bookkeeping.

\section{Safe collapse: the must-quotient}
\label{sec:quotient}

The feasible family is a set of equivalence relations, and the
practitioner wants \emph{one} deterministic output.  The only generally
safe deterministic collapse is the intersection of the feasible
relations: define $\Must_D(c,d)$ to hold when every feasible hypothesis
merges $c$ and $d$.  The mechanization proves this intersection is
itself an equivalence relation (\path{mustSame_equivalence}), packaged
as the setoid \path{mustSetoid}, restricted to active claims as
\path{activeMustSetoid}, with quotient type \path{CanonicalFrame}.
Collapsing by $\Must$ therefore never invents an identification: it
reports exactly the merges that survive in \emph{all} feasible worlds.

The remaining pairs are not noise; they are classified.  For a
\emph{consistent} corpus every pair is must-merged, cannot-merged, or
explicitly uncertain (\path{coreference_exhaustive}); the deterministic
classifier \path{corefStatus} returns
$\{\mathrm{must},\mathrm{cannot},\mathrm{uncertain},
\mathrm{inconsistent}\}$, with the inconsistent case kept separate
precisely because universal quantification over an empty feasible family
is vacuous (\path{must_and_cannot_implies_inconsistent}).  The statuses
are stable in the useful direction: settled statuses persist under
corpus growth (\path{mustSame_mono}, \path{cannotSame_mono}), while
uncertainty can only be resolved, never created, by adding evidence
(\path{uncertain_antitone}).

Why not collapse more aggressively --- say, quotient by ``some feasible
world merges them''?  Because the union of equivalence relations across
feasible worlds is \emph{not} an equivalence relation.  The finite
countermodel \path{Ste.DynamicFrame.Counterexample} exhibits three
claims $a,b,c$ and two feasible worlds (one merging $\{a,b\}$, one
merging $\{b,c\}$) such that $\May(a,b)$ and $\May(b,c)$ hold while
$\Cannot(a,c)$ holds: \path{maySame_not_transitive}.  The uncertainty
overlay fails transitivity the same way
(\path{uncertain_not_transitive}).  Any pipeline that
transitively closes may-links or uncertain-links manufactures a merge
($a\sim c$) that every feasible hypothesis refutes.

\begin{remark}[Qualification]
The must-quotient is safe \emph{on a consistent corpus}.  If the
feasible family is empty, $\Must$ holds vacuously for every pair and the
quotient degenerates; the mechanization guards this by carrying the
consistency premise in \path{coreference_exhaustive} and by routing
empty families to the separate \path{corefStatus} value
\texttt{inconsistent}.  Check feasibility before collapsing.
\end{remark}

\emph{When to reach for it.}  Whenever a downstream consumer demands a
single partition.  Emit the must-quotient plus the uncertain-pair list;
never a transitive closure of possibilities.

\section{Certification: minimal supports}
\label{sec:support}

A settled conclusion ($\Must$ or $\Cannot$ on some finite corpus) drags
the whole corpus behind it as justification.  That justification can be
shrunk.  Define \path{SupportsMust}~$(D,c,d)$: the finite corpus $D$ is
consistent and forces $c,d$ together (analogously
\path{SupportsCannot}).  The mechanization proves that every finitely
supported conclusion has an \emph{inclusion-minimal} finite subcorpus
that still supports it: \path{exists_minimal_mustSupport} and
\path{exists_minimal_cannotSupport}, instances of the general
finite-minimization principle \path{exists_minimal_subfinset}.  The
minimal support is a certificate: a smallest-by-inclusion set of
documents such that removing any one breaks the conclusion.

The mechanization is deliberately semantic here: the theorems assert
existence of minimal supports, not that finding one is cheap.  Note also
that consistency is part of the support predicate --- a subcorpus that
forces the pair only vacuously (by being inconsistent) does not count as
support.

\emph{When to reach for it.}  Explanation, audit, and caching.  A
minimal support is simultaneously the provenance answer (``why are these
merged?''), the invalidation key for dynamic updates (the conclusion
survives any deletion disjoint from its support), and a compressed proof
obligation for human review.

\section{Composition: gluing local families with an exact obstruction}
\label{sec:gluing}

Large problems are attacked in pieces: compute feasible families on
subcorpora $D$ and $E$, then combine.  The obstruction to combining is
exactly the genuinely cross-piece constraints.  For an ambient
hypothesis $h$, \path{crossObstruction} collects the constraints whose
support lies in $D\cup E$ but in neither piece alone and which $h$
violates; the classification theorem \path{mem_feasible_union_iff}
proves
\[
h \in \mathrm{feasible}(D\cup E)
\iff
h \in \mathrm{feasible}(D)\ \wedge\ h \in \mathrm{feasible}(E)\ \wedge\
\mathcal O_{D,E}(h)=\emptyset .
\]
Two local world views are \path{Compatible} when they expose the same
ambient hypothesis; a \path{GlobalExtension} is a world on $D\cup E$
restricting to both.  Existence is settled exactly:
\path{globalExtension_nonempty_iff} proves a global extension exists iff
the views are compatible and the cross obstruction vanishes, with the
witness constructed by \path{glue}.  Uniqueness is a theorem, not a
hope: the extension type is a subsingleton
(\path{globalExtension_subsingleton}) --- at most one global world
extends a fixed pair of local worlds.

\begin{remark}[Qualification]
Uniqueness is relative to the compatibility notion mechanized here:
all views live in one ambient hypothesis type, and compatibility means
literal equality of the underlying hypothesis.  In a setting where local
views are related by a looser correspondence, uniqueness would have to
be re-proven for that correspondence.  The obstruction set is
computable in the finite case, but it is an obstruction \emph{per
ambient hypothesis}: gluing whole families means discharging it per
surviving candidate.
\end{remark}

\emph{When to reach for it.}  Divide-and-conquer over provenance:
shard the corpus, shrink each shard's family locally (cheap), and pay
the cross-constraint check only at the join.  The theorem tells you
exactly which constraints can obstruct --- the straddlers --- so the
join cost is proportional to genuine coupling, not to corpus size.

\section{The wall: \texorpdfstring{$2^N$}{2\^{}N} extensional states, and why it is tight}
\label{sec:wall}

How far can generic shrinking go?  The finite mechanization
\path{Ste.DynamicFrame.FiniteSolver} makes the ceiling exact.  When the
hypothesis universe is finite with $N=$ \path{maxWorlds} ambient worlds
and satisfaction is decidable, the solver state is represented exactly
by a finite set: filtering implements the meet
(\path{coe_applyFinite}), and the executable run denotes precisely the
abstract STE run (\path{coe_runFinite}, \path{coe_solveFinite}) ---
partial feasibility on the nose.  Total reduction is bounded: a schedule
can contain at most $N$ properly shrinking steps
(\path{reductionCost_le_maxWorlds}).  The state space is
$2^N$ exactly (\path{card_allExactStates}).

The bound is tight, and this is the pivot of the whole note: a single
crisp property set can select \emph{any} chosen subset of the universe
(\path{arbitrary_subset_as_single_property}).  Unrestricted crisp STE
therefore realizes all $2^N$ extensional states, and no uniform
sub-exponential family of states can be complete for it.  No amount of
clever generic machinery --- meets, restriction, quotients, supports,
gluing --- improves the worst case, because the worst case is
realizable by the inputs themselves.  Generic reduction stops here.
Everything further must buy its compression with \emph{structure}: an
assumption on the shape of the generated property sets.

\section{Structural reduction: width-bounded compression}
\label{sec:structure}

Structure enters through \emph{scopes}: a property set over a product
space $\prod_v A_v$ has support $\sigma$ when it depends only on the
variables in $\sigma$ (\path{HasSupport}).  Scoped constraints induce a
constraint (hyper)graph, and when that graph is tree-like the feasible
family compresses.  The mechanization builds the full chain from scope
accounting to a polynomial-size faithful representation.

\subsection{Scope accounting and cut elimination}

Joining constraints unions their scopes (\path{HasSupport.inter_union},
family form \path{HasSupport.iInter}); conditioning (fixing a variable)
removes it from the scope and commutes with joins
(\path{condition_inter}).  The pivotal dissolution step is
\path{cut_elimination}: if two constraints share only a cut variable
$v$, conditioning on $v$ leaves constraints with \emph{disjoint} scopes
--- the problem falls apart into independent blocks, and joint
feasibility becomes independent per-block feasibility
(\path{condition_inter_nonempty_iff}).

The space currency is the \path{table}: the projection of a constraint
to its scope.  A supported constraint is faithfully recovered from its
table (\path{HasSupport.eq_preimage_table}), and the table of a bag of
at most $w+1$ variables over alphabets of size at most $k$ has at most
$k^{w+1}$ rows (\path{table_encard_le}, \path{table_encard_le_pow}).

\subsection{Elimination and adaptive consistency}

Bucket elimination \citep{dechter1999bucket,dechter2003constraint} is
mechanized as a fold of conditioning steps.  One step joins a bucket and
eliminates its variable, staying within the bag bound
(\path{elimination_step}); a whole order of width $w$ with $n$ steps
materializes at most $n\cdot k^{w+1}$ total table rows
(\path{elimination_order_total_bound},
\path{elimination_order_table_total_bound},
\path{bucketEliminate_total_space}).  A \emph{complete} order --- one
covering every scope --- decides the instance: the residue is
\path{Set.univ} or $\emptyset$ (\path{eliminate_eq_univ_or_empty},
bucket form \path{bucketEliminate_decides}), i.e.\ the feasibility query
is answered without ever materializing the feasible family.

The fused solver statement is
\path{projectBucketEliminate_complete_solver}: for any complete order of
width $w$, one theorem simultaneously decides feasibility and
infeasibility, extracts a concrete global solution backtrack-free when
feasible, and certifies the width bound
\path{inducedTreewidth} $\le w$.  Adaptive consistency proper --- the
decreasing-order run whose separators are the algorithm's own recorded
message scopes --- is \path{decreasingRun_adaptiveSolver}.

\subsection{Junction trees: a faithful polynomial representation}

Elimination's trace can be kept as a data structure: the
\path{junctionTables} (one bag table per materialized bucket) with total
size \path{junctionSize}.  This representation is \emph{faithful}: a
scoped constraint is exactly the preimage of its bag table, and a global
assignment is feasible for the whole instance iff it restricts into
every bag table (\path{mem_joinConstraint_iff_restrict},
\path{junctionTables_faithful}).  The size theorem
\path{junctionTree_size_le} bounds the representation by
$n\cdot k^{\,\mathrm{tw}+1}$ over the induced treewidth; at fixed width
the bound is linear in the number of variables
(\path{junctionTree_size_linear}, existence form
\path{junctionTree_size_le_card} at the minimal duplicate-free width
\path{junctionWidth}, which dominates the induced treewidth by
\path{inducedTreewidth_le_junctionWidth} --- that the two coincide is
not mechanized and remains outlook).  The upshot: a feasible family of
size exponential in $n$ is represented, queried for feasibility, and
sampled from, in $O(n)$ space at bounded width.

\subsection{Factorization over partitions}

When the constraint graph is disconnected the compression is exact and
multiplicative.  For blocks pulled back along coordinate projections
(\path{pullbackFamily}), the feasible family is exactly the product of
per-block families (\path{feasibilitySet_pullbackFamily}) and its
cardinality is the product of per-block cardinalities
(\path{encard_feasibilitySet_blocks}; two-level version
\path{encard_feasibilitySet_blockFamilies}; domination bound
\path{encard_feasibilitySet_blocks_le}).  An estimator pays
$\sum_i \mathrm{cost}(C_i)$ across blocks instead of
$\prod_i \mathrm{cost}(C_i)$ for the joint family: the exponential is
confined inside each block.

\subsection{Representation bounds: how compressible is a family?}

The rectangle-cover machinery of \path{Ste.RepresentationBounds}
quantifies compressibility in the communication-complexity sense
\citep{kushilevitz1997communication}.  The cover number
\path{rectCoverNumber} --- the least number of variable-separable
product sets whose union is the family --- is sandwiched by machine
checked bounds: it is at most the family's cardinality
(\path{rectCoverNumber_le_encard}), at least every fooling set
(\path{foolingNumber_le_rectCoverNumber}), and equals $1$ exactly when
the \v{C}ech obstruction of the singleton cover vanishes
(\path{rectCoverNumber_eq_one_iff_cechVanishes}).  On the maximally
coupled family the cover number is computed exactly
(\path{rectCoverNumber_allEqual}).  These give the practitioner a
diagnostic: a small fooling set in the shrunken family is a certificate
that no small rectangular (per-variable factored) representation
exists.

\subsection{The Carlson bridge: reducing the problem class itself}

Structural reduction can act one level up, on the hypothesis universe
before any property set is applied.  Carlson's dissertation recasts
decryption as STE \citep{carlson2012}: keys are hypotheses, each
intercepted ciphertext induces a property set
(\path{keyPropertySet}), and the feasible key set is literally an STE
feasibility set (\path{feasibleKeys}, \path{mem_feasibleKeys}), with
more traffic only narrowing it (\path{feasibleKeys_antitone}).  The
residual family size after one observation is counted exactly:
$(|A|-|T|)!$ keys remain consistent, where $T$ is the observed symbol
set (\path{card_consistent_keys}, Carlson's Lemma~4.1), out of $|A|!$
substitution keys (\path{card_substitution_keys}); identification is
the ideal case \path{Unicity}.

The reduction results then collapse cipher \emph{classes}: every
permutation-cipher key is a substitution key via the injective
realization \path{coperm}
(\path{permutation_reduces_to_substitution}, sharpened by
\path{coperm_injective}), the reverse reduction is impossible for
$n\ge 2$ --- there are strictly more substitution keys than
permutation keys, so no faithful embedding is onto
(\path{card_permutation_lt_card_substitution},
\path{substitution_not_reducible_to_permutation}) --- and any
permutation--substitution--permutation composite is again a single
substitution key (\path{psp_reduces_to_substitution}); cf.\ the same
reduction in the applied literature
\citep{ghosh2021isomorphic,carlson2022equivalence}.  For the STE
practitioner this is the template for \emph{class-level} shrinking:
prove one hypothesis class embeds in another, attack the larger class
once, and inherit the feasible-family analysis for every problem in the
smaller class --- while the counting inequality warns which direction
the reduction cannot go.

\section{The lower-bound wall for structure}
\label{sec:lowerbound}

Structure is not optional decoration; when it is absent, blow-up is
forced.  The singleton-cover gluing characterization
\path{cechVanishes_iff_rectangular} proves that margin-based local
reasoning is exact --- every compatible family of per-variable local
sections glues --- \emph{iff} the family is a rectangle: no coupling, no
obstruction.  The $n$-fold coupling \path{allEqual} (all coordinates
equal over an alphabet $\alpha$) sits at the opposite extreme: every
per-variable margin is full, so all $|\alpha|^n$ local-section families
are compatible (\path{compatibleFamilies_allEqual_encard}), yet only the
$|\alpha|$ constant vectors are actually feasible
(\path{allEqual_encard}).  The compatible-but-infeasible gap --- the
phantom assignments any per-variable procedure must entertain and the
coupling then kills --- is exactly $|\alpha|^n-|\alpha|$
(\path{cechObstruction_allEqual}), hence at least $2^n-2$ once
$|\alpha|\ge 2$ (\path{cechObstruction_allEqual_ge}): exponential in the
number of variables while the feasible family itself has constant size.
For $n\ge 2$ the constraint is genuinely non-rectangular and the
obstruction genuinely nonvanishing (\path{allEqual_not_rectangular},
\path{allEqual_not_cechVanishes}).

\begin{remark}[Honest boundary]
What is mechanized is exactly: the iff for the singleton cover, and the
exact and asymptotic obstruction counts for \path{allEqual}.  Not
mechanized (and not claimed): that \path{allEqual}'s induced width is
minimal in the formal \path{Ste.Treewidth} sense, and the fully general
statement that \emph{every} width-$w$ representation scheme for
\path{allEqual} needs size $2^{\Omega(n)}$; the cohomological picture
for covers with overlaps \citep{abramsky2011sheaf} remains outlook.
The practical reading is nevertheless sharp: before investing in
elimination orders and junction trees, measure the coupling; a
complete-graph-like instance pays the exponential somewhere no matter
how the pipeline is arranged.
\end{remark}

\section{Practitioner contract}
\label{sec:contract}

The techniques compose into a fixed order of operations for an STE
attack.

\begin{enumerate}
\item \textbf{Generate} the property-set family from the information
  sources; record provenance supports.
\item \textbf{Apply} audited property sets with the exact meet
  (\path{applyConstraint}, \path{solve_eq_feasible}); schedule freely
  (idempotent, commutative); log the exact deltas
  (\path{apply_union_reduction}).
\item \textbf{Restrict / garbage-collect} across instance changes:
  reuse the family under corpus growth via the delta law
  (\path{mem_feasible_extension}); drop retired claims by partial
  restriction (\path{claimRestriction_exists_iff}); never invert a
  merge by bookkeeping (retraction asymmetry).
\item \textbf{Check feasibility}, and treat emptiness as detection of
  invalid information
  (\path{exists_unfair_of_feasibilitySet_eq_empty}).
\item \textbf{Collapse safely} by the must-quotient only
  (\path{mustSame_equivalence}, \path{CanonicalFrame}); publish
  must/cannot/uncertain statuses (\path{corefStatus}); never
  transitively close may-links (\path{maySame_not_transitive}).
\item \textbf{Certify} settled conclusions by inclusion-minimal
  supports (\path{exists_minimal_mustSupport}).
\item \textbf{Compose} shard-wise with the exact gluing obstruction
  (\path{globalExtension_nonempty_iff}); uniqueness of extension is
  guaranteed under compatibility
  (\path{globalExtension_subsingleton}).
\item \textbf{Compress structurally} when scopes permit: bucket
  elimination and adaptive consistency decide and solve within induced
  width (\path{projectBucketEliminate_complete_solver}); keep the
  junction-tree tables as the $O(n)$ faithful representation at fixed
  width (\path{junctionTree_size_le}); factor across disconnected
  blocks (\path{encard_feasibilitySet_blocks}); reduce hypothesis
  classes when an embedding is proven
  (\path{permutation_reduces_to_substitution}).
\item \textbf{Know the walls}: generically, $2^N$ extensional states
  are realizable and tight
  (\path{card_allExactStates},
  \path{arbitrary_subset_as_single_property}); structurally, maximal
  coupling forces an obstruction of at least $2^n-2$
  (\path{cechObstruction_allEqual_ge}).
\end{enumerate}

\medskip
\noindent
Finally, the preservation ledger --- which queries each step provably
preserves:

\begin{center}
\begin{tabular}{@{}lll@{}}
\toprule
Step & Preserves & Authority \\
\midrule
Exact meet & all queries (it \emph{defines} them) &
  \path{solve_eq_feasible} \\
Restriction ($D\subseteq E$) & must/cannot upward; may/uncertain
  downward &
  \path{mustSame_mono}, \path{maySame_antitone} \\
Must-quotient & must; statuses of remaining pairs &
  \path{coreference_exhaustive} \\
Minimal support & the certified must/cannot fact &
  \path{exists_minimal_mustSupport} \\
Gluing & feasibility on the union, exactly &
  \path{mem_feasible_union_iff} \\
Elimination (complete order) & feasibility + one witness &
  \path{projectBucketEliminate_complete_solver} \\
Junction tables & full membership (hence all queries) &
  \path{mem_joinConstraint_iff_restrict} \\
Factorization & full family, as a product &
  \path{feasibilitySet_pullbackFamily} \\
\bottomrule
\end{tabular}
\end{center}

\medskip
\noindent
Restriction deserves the final caution, because its ledger row is
asymmetric by theorem, not by implementation: settled must/cannot facts
persist as the corpus grows (\path{mustSame_mono},
\path{cannotSame_mono}) but may be lost under deletion, while may/
uncertain facts transfer in the opposite direction only
(\path{maySame_antitone}, \path{maySeparate_antitone},
\path{uncertain_antitone}).  A pipeline that caches query answers across
corpus edits must cache them with the direction of the edit --- which
is, in the end, the whole discipline this note proposes: shrink as hard
as the theorems allow, and no harder.

\bibliographystyle{plainnat}
\bibliography{../../refs}

\end{document}
